\begin{table}[t]
\centering
\caption{Summary Statistics: Pre-Treatment Industry Characteristics (2012--2013)}
\label{tab:summary}
\footnotesize
\begin{tabular}{llrrrrr}
\toprule
NAICS & Industry & Mean Emp & SD Emp & Mean Estab & Mean Payroll & County- \\
 & & & & & (\$000s) & Years \\
\midrule
\multicolumn{7}{l}{\textit{Panel A: Treated Industries (Software-Intensive)}} \\
334 & Computer/Electronics & 2,266 & 4,976 & 36.5 & 180.0 & 562 \\
511 & Publishing/Software & 1,404 & 4,203 & 40.3 & 143.3 & 988 \\
518 & Data Processing/Internet & 1,481 & 2,495 & 45.5 & 136.8 & 463 \\
\midrule
\multicolumn{7}{l}{\textit{Panel B: Control Industries (Not Software Patent-Dependent)}} \\
325 & Chemicals/Pharma & 1,101 & 1,994 & 20.9 & 79.8 & 886 \\
336 & Transportation Equip & 1,895 & 4,404 & 19.1 & 118.0 & 715 \\
339 & Misc Manufacturing & 768 & 1,794 & 42.0 & 41.1 & 1,057 \\
\bottomrule
\end{tabular}
\begin{tablenotes}
\item \textit{Notes:} Data from Census County Business Patterns (CBP), 2012--2013.
Treated industries are those whose core activities rely on software patents affected by
\textit{Alice Corp. v. CLS Bank} (2014). Control industries produce physical goods or
pharmaceuticals whose patents are not subject to Section 101 eligibility challenges.
Employment and payroll are annual county-level values.
\end{tablenotes}
\end{table}
